
\begin{ledgroup}
\footnotesize
\pstart
\noindent
Mit N.\;\ref{57034} und N.\;\ref{57035} liegen zwei eigenhändig datierte "\textit{Schedae}" vor, die sich dem statischen Thema aufeinander wirkender Hebel und der Bestimmung der darauf folgenden Bewegungen widmen. 
Solche Hebel bieten auch ein gutes Modell für die Wechselwirkung der ineinander verschränkten Zähne zweier Zahnräder, deren Eigenschaften Leibniz in N.\;\ref{intro_vectibus-conjugatis} ebenfalls behandelt ("de figuris dentium in rotis").
\pend
\pstart
Leibniz hatte sich bereits in der Pariser Zeit mit diesen Fragen beschäftigt: In unmittelbar benachbarten Signaturen (LH~XXXVII~3, Bl.~77\textendash80) sind vier Konzepte \cite{02160}"De vectibus conjugatis" aus der zweiten Hälfte von 1674 überliefert (ediert in \textit{LSB} VIII,~2 N.\;45). 
Die Diagramme \bracks{\textit{Fig.~1}}, \bracks{\textit{Fig.~3}} und \bracks{\textit{Fig.~8}} von N.\;\ref{57034} greifen bspw. die Figuren des ersten Pariser Konzepts (\textit{LSB} VIII,~2 N.\;45\textsubscript{1}) auf.
Leibnizens Angabe "resumtae priores cogitationes" in der Überschrift von N.\;\ref{57034} bezieht sich daher aller Wahrscheinlichkeit nach auf die vier Texte von 1674.
Ein weiterer Zeuge seiner Beschäftigung mit dem Thema liegt mit der Notiz "De curvis inveniendis ope vectium conjugatorum" vor, die gemäß ihrer wahrscheinlichen Entstehungszeit (Dezember 1674 bis 1675) in diesem Band als Nachtrag ediert wird (N.\;\ref{Cc2_847-2}).
\pend
\pstart
Der Text der zweiten Konzepts N.\;\ref{57035} ist ursprünglich als direkte Fortsetzung des ersten entstanden, wie der Kustos "Longe secus est" am Ende von N.\;\ref{57034} (Bl.~76~v\textsuperscript{o}) bezeugt.
Erst nachträglich hat Leibniz die Einteilung in "\textit{Scheda prima}" und "\textit{secunda}" vorgenommen, 
indem er auf Bl.~75~r\textsuperscript{o}, hier \refpassage{37_03_075-076_1a}{37_03_075-076_1b}, und Bl.~81~r\textsuperscript{o}, hier \refnote{FnB:37_03_081_scheda-2da}, entsprechende Überschriften hinzufügte.
\pend
\end{ledgroup}
